% §1.4 Sampling and Estimating Rates
%
% Cross-reference resolution status:
%
% ALREADY DEFINED (within chapter 1, by this section):
%   \label{sec:sampling-rates}              -- this section
%
% RESOLVED:
%   \Cref{sec:quantifying-noise}            -- §1.2
%   \Cref{tab:contingency}                  -- §1.2
%   \Cref{ch:classification-measures}       -- chapter 7
%   \Cref{sec:axiom-1}                      -- chapter 4 §4.2 (independence)
%
% FORWARD REFS to later chapter 1 sections:
%   \Cref{sec:why-atom}                     -- §1.5
%
% Monograph register: states the binomial estimator, its variance, and the
% Wald interval without derivation (full theory in chapter 7), frames the
% notebook as a reproducibility artifact, and salvages the former
% correlated-errors exercise as a remark (independence vs. stability).

\section{Sampling and Estimating Rates}
\label{sec:sampling-rates}

In practice the rates are not given: a classifier arrives without a label
stating its false-positive rate, and a channel must be characterized from
observed traffic. We state how to estimate $\fprate$ and $\fnrate$ from a
finite sample and how to quantify the uncertainty. The full theory of these
estimators (bias, variance, efficiency, and interval forms) is developed in
\Cref{ch:classification-measures}.

\subsection*{Binomial estimators}

Run a noisy bit $n$ times, with $n_0$ trials of truth $\bot$ and $n_1 = n -
n_0$ of truth $\top$. Each $\bot$-trial reports $\top$ independently with
probability $\fprate$, so the false-positive count is binomial,
$\mathrm{FP} \sim \mathrm{Binomial}(n_0, \fprate)$, and likewise
$\mathrm{FN} \sim \mathrm{Binomial}(n_1, \fnrate)$. The row-conditional
proportions
\[
  \hat{\fprate} = \frac{\mathrm{FP}}{n_0},
  \qquad
  \hat{\fnrate} = \frac{\mathrm{FN}}{n_1}
\]
are unbiased, $\mathbb{E}[\hat{\fprate}] = \fprate$, with variances
\[
  \mathrm{Var}(\hat{\fprate}) = \frac{\fprate(1 - \fprate)}{n_0},
  \qquad
  \mathrm{Var}(\hat{\fnrate}) = \frac{\fnrate(1 - \fnrate)}{n_1}.
\]
The variance shrinks as $1/n_0$, and its numerator $\fprate(1 - \fprate)$ peaks
at $\fprate = 1/2$, where the noisy bit is hardest to characterize, and
vanishes at $\fprate \in \{0, 1\}$, where one trial settles the matter.

\subsection*{A confidence interval}

For large $n_0$ the central limit theorem makes $\hat{\fprate}$ approximately
Normal with mean $\fprate$ and variance $\fprate(1 - \fprate)/n_0$. Substituting
the estimate into the variance gives the Wald interval at level $1 - \alpha$,
\[
  \hat{\fprate} \;\pm\; z_{\alpha/2}
    \sqrt{\frac{\hat{\fprate}(1 - \hat{\fprate})}{n_0}},
\]
with $z_{\alpha/2}$ the upper-$\alpha/2$ standard-Normal quantile. The Wald form
is asymptotically correct but under-covers at small $n_0$ when $\fprate$ is near
$0$ or $1$, and can even leave $[0,1]$; the Wilson, Clopper-Pearson, and
Jeffreys intervals handle the boundary better, and \Cref{ch:classification-measures}
develops them.

\subsection*{Reproducibility}

The companion notebook \texttt{bernoulli-py/notebooks/ch01-noisy-bit.ipynb}
reproduces these claims: it estimates $\fprate$ from a simulated binary
symmetric channel across sample sizes spanning several orders of magnitude,
exhibiting the $1/\sqrt{n_0}$ concentration of $\hat{\fprate}$ about the true
rate; and it measures the empirical coverage of the $95\%$ Wald interval,
showing it approach nominal at moderate $n_0$ and fall below nominal as $n_0$
shrinks or $\fprate$ approaches a boundary, the under-coverage noted above.

\begin{remark}[Independence is separate from rate-stability]
\label{rem:independence-vs-stability}
The binomial model assumes the per-query error events are independent.
Independence is a distinct assumption from the stability of the per-input
rates: a system can have stable rates $\fprate$ and $\fnrate$ yet errors that
are correlated across queries. The point estimates $\hat{\fprate},
\hat{\fnrate}$ remain valid in that case, but the variances above, and every
confidence interval and sample-size requirement built on them, misstate the
true uncertainty. Independence is named as Axiom~1 in
\Cref{sec:axiom-1}, and chapter~16 takes correlated-error systems as an
explicit subject.
\end{remark}
